\begin{table}[H]
\centering
\caption{Κύριες προκλήσεις του \ENG{IEEE} \ENG{LOM}}
\label{tab:ieee_lom_challenges}
\small
\setlength{\tabcolsep}{5pt}
\renewcommand{\arraystretch}{1.25}
\begin{chaptertablebox}
\begin{tabular}{|p{3.5cm}|p{9.5cm}|}
\hline
\textbf{Κατηγορία} & \textbf{Προβλήματα} \\
\hline
\hline
\textbf{Πολυπλοκότητα} & 76 στοιχεία - υπερβολική πολυπλοκότητα για \ENG{manual} \ENG{creation}.\newline
Δυσκολία κατανόησης και συμπλήρωσης όλων των πεδίων.\newline
Υψηλό κόστος παραγωγής \ENG{metadata}. \\
\hline
\textbf{Αμφισημία} & Πολλά πεδία έχουν ασαφείς ορισμούς.\newline
Διαφορετικές ερμηνείες από διαφορετικούς \ENG{implementers}.\newline
Έλλειψη συνέπειας μεταξύ \ENG{repositories}. \\
\hline
\textbf{Χαμηλή \ENG{Adoption} \ENG{Rate}} & Οι περισσότεροι \ENG{content} \ENG{creators} δεν συμπλήρωναν πλήρη \ENG{LOM} \ENG{metadata}.\newline
Τα περισσότερα αποθετήρια χρησιμοποιούσαν μόνο μικρό υποσύνολο πεδίων.\newline
\ENG{Automated} \ENG{metadata} \ENG{generation} ήταν περιορισμένη. \\
\hline
\textbf{Ελλιπής Υποστήριξη Εργαλείων} & Περιορισμένα \ENG{authoring} \ENG{tools} με ενσωματωμένη υποστήριξη \ENG{LOM}.\newline
Δύσχρηστα \ENG{metadata} \ENG{editors}.\newline
Περιορισμένη ενσωμάτωση με \ENG{content} \ENG{development} \ENG{workflows}. \\
\hline
\end{tabular}
\end{chaptertablebox}
\end{table}
